% Vietnamese translation for OI Public Library
% Source-File: IOI中国国家候选队论文/2009/曹钦翔/从“k倍动态减法游戏”出发探究一类组合游戏问题.ppt
\documentclass[11pt]{article}
\usepackage{oipl-vn}

\newcommand{\mex}{\operatorname{mex}}

\title{Từ trò chơi trừ động theo bội \(k\) đến một lớp bài toán trò chơi tổ hợp\\{\large Ghi chú theo slide}}
\author{Cao Qinxiang\\Trường Trung học số 2 trực thuộc Đại học Sư phạm Hoa Đông}
\date{2009}

\begin{document}
\maketitle

\origtitle{From k-dynamic subtraction games to combinatorial games}
\sourcepath{IOI-China-Candidate-Team-Papers/2009/Cao-Qinxiang/k-dynamic-subtraction-games-slides.ppt}

\section{Mạch trình bày}

Slide bắt đầu từ trò chơi trừ động theo bội \(k\), trong đó số lượng lấy ở
lượt hiện tại bị ràng buộc bởi lượt trước. Từ ví dụ này, tác giả mở rộng sang
cách phân tích trò chơi tổ hợp bằng trạng thái N/P, hàm SG, tổng Nim và một
số kỹ thuật không thuần quy hoạch động.

\section{Các điểm kỹ thuật}

\begin{itemize}
  \item Trạng thái P là trạng thái mọi nước đi đều dẫn đến N; trạng thái N có
  ít nhất một nước đi đến P.
  \item Quy hoạch động tổng quát có thể tính thắng thua, nhưng số nước đi lớn
  khiến cần khai thác đơn điệu hoặc duy trì tập quyết định tốt.
  \item Hàm SG dùng \(\mex\) của các giá trị con; tổng của nhiều trò độc lập
  được xử lý bằng xor.
  \item Một số bài có quy luật đối xứng hoặc chiến lược tham lam, giúp tránh
  liệt kê toàn bộ đồ thị trạng thái.
\end{itemize}

\section{Ghi chú}

Thông điệp thực dụng của slide là không nên dừng ở công thức SG tổng quát.
Sau khi mô hình hóa trạng thái, hãy tìm tính đơn điệu, cấu trúc chu kỳ, đối
xứng hoặc cách duy trì nước thắng để giảm chi phí tính toán.

\end{document}
